\section{Walsh transforms and moment lifts}\label{sec:preliminaries}

We first introduce the operator norms that measure Walsh decay, then derive
the moment lifts needed to compare those norms at different dimensions.
Throughout, Banach spaces are real. All unspecified averages on finite
sets are uniform, and distinct averaging variables are independent.
The notation \(B_Y\) means the closed unit ball of \(Y\).
For a finite probability space \(\Omega\), the space \(L_p(\Omega;Y)\)
has its probability-normalized norm. By contrast,
\(\ell_2^m(Y)\) has norm \((\sum_{i=1}^m\|y_i\|^2)^{1/2}\).
Duality on finite probability spaces uses the expectation pairing;
in particular,
\[
\begin{gathered}
L_2(\Omega;Y)^*=L_2(\Omega;Y^*),\\
L_1(\Omega;Y)^*=L_\infty(\Omega;Y^*),\qquad
L_\infty(\Omega;Y)^*=L_1(\Omega;Y^*).
\end{gathered}
\]

\subsection{Walsh norms}

Let \(V_n=\F_2^n\), with its standard dot pairing, and put
\(\chi(u,v)=(-1)^{u\cdot v}\). For \(T:G\to F\), define
\[
(W_n^Tx)(v)=\E_{u\in V_n}\chi(u,v)Tx(u),\qquad
w_n(T)=\|W_n^T:L_2(V_n;G)\to L_2(V_n;F)\|.
\]
We allow \(n=0\), so \(w_0(T)=\|T\|\).
Finite-sum duality gives the equivalent formula, with
\(x_u\in G\) and \(y_v\in F^*\),
\begin{equation}\label{eq:walsh-bilinear}
w_n(T)=\sup\left|
\E_{u,v\in V_n}\chi(u,v)y_v(Tx_u)\right|,
\qquad
\E_u\|x_u\|^2\le1,\quad \E_v\|y_v\|^2\le1.
\end{equation}

\begin{lemma}[Basic Walsh estimates]\label{lem:walsh-basic}
For every bounded \(T:G\to F\):
\begin{enumerate}[label=\textup{(\roman*)}]
\item \(2^{-n/2}\|T\|\le w_n(T)\le\|T\|\), and \(w_n(T)\) is
nonincreasing in \(n\).
\item If \(T_\Omega:L_2(\Omega;G)\to L_2(\Omega;F)\) acts pointwise
as \(T\), where \(\Omega\) is finite, then
\(w_n(T_\Omega)=w_n(T)\).
\item If \(T\) has finite rank, there is \(D_T<\infty\) such that
\begin{equation}\label{eq:rank-dependent-decay}
w_n(T)\le D_T2^{-n/2}\qquad(n\ge0).
\end{equation}
\end{enumerate}
Moreover, for \(A,B:G\to F\),
\begin{equation}\label{eq:walsh-continuity}
|w_n(A)-w_n(B)|\le\|A-B\|.
\end{equation}
\end{lemma}
\begin{proof}
The upper bound follows from the triangle inequality and
Cauchy--Schwarz; the lower bound follows by taking input supported
at one \(u\), scaled by \(2^{n/2}\). On \(V_{n+j}=V_n\oplus V_j\),
first apply \(W_n^T\) in the first coordinate and then the normalized
Walsh transform with the identity of \(F\) in the second coordinate.
The latter has norm at most \(1\). Averaging the squared estimates
in the unused coordinates proves monotonicity.
Interchanging finite \(L_2\) coordinates proves (ii), with constant
functions on \(\Omega\) giving the reverse inequality.
For (iii), factor a finite-rank \(T\) as \(G\to H\to F\), with \(H\)
a finite-dimensional Hilbert space. Orthogonality of the Walsh
matrix gives \(\|W_n^{\Id_H}\|=2^{-n/2}\), proving (iii).
Finally, \(W_n^A-W_n^B=W_n^{A-B}\), so the upper bound in (i)
and the reverse triangle inequality prove \eqref{eq:walsh-continuity}.
\end{proof}

\subsection{Fixed-length cotype estimates}

Define \(c_m(Y)\) to be the least constant in
\begin{equation}\label{eq:fixed-cotype}
\left(\sum_{i=1}^m\|y_i\|^2\right)^{1/2}
\le c_m(Y)\left(\E_\eps\left\|\sum_{i=1}^m\eps_i y_i\right\|^2\right)^{1/2}.
\end{equation}
Thus, if \(Y\) has cotype \(q\) with constant \(C_q(Y)\), then
\begin{equation}\label{eq:fixed-from-cotype}
c_m(Y)\le C_q(Y)m^{1/2-1/q}.
\end{equation}
Squaring \eqref{eq:fixed-cotype} and averaging proves
\begin{equation}\label{eq:cotype-amplification}
c_m(L_2(\Omega;Y))\le c_m(Y).
\end{equation}
The same observation applies to duals of these finite \(L_2\) sums.

The following estimate is the \(L_1\)-to-\(L_2\) case of Kahane's
inequality; see \cite[Theorem~8.1]{Pis11IHP}.
We include a stopping-time proof with an explicit, sufficient constant
so that the subsequent moment lifts are self-contained.

\begin{lemma}[Moment comparison]\label{lem:kahane}
For every finite Rademacher sum \(S=\sum_{i=1}^m\eps_i z_i\) in a Banach
space,
\begin{equation}\label{eq:kahane}
(\E\|S\|^2)^{1/2}\le L\,\E\|S\|,\qquad L=24.
\end{equation}
\end{lemma}
\begin{proof}
Put \(M=\E\|S\|\) and \(b=\max_i\|z_i\|\).
Since \(z_i=\E(\eps_iS)\), we have \(b\le M\).
The case \(M=0\) is immediate.

Let \(S_j=\sum_{i\le j}\eps_i z_i\), and stop at the first \(j\) with
\(\|S_j\|>u\). Conditional on a stopped prefix \(z\), the remaining
sum is symmetric; at least one of \(\|z+h\|\) and \(\|z-h\|\)
exceeds \(u\). Hence
\[
\mathbb P(\text{a stop occurs})\le2\mathbb P(\|S\|>u).
\]
For any fixed \(j\), symmetry of the independent prefix similarly gives
\[
\mathbb P\left(\left\|\sum_{i>j}\eps_i z_i\right\|>v\right)
\le2\mathbb P(\|S\|>v).
\]
At the first stop, \(\|S_j\|\le u+b\). Decomposing according to its
time and using independence of the remaining sum therefore yields
\[
\mathbb P(\|S\|>u+b+v)
\le4\mathbb P(\|S\|>u)\mathbb P(\|S\|>v).
\]
Take \(u=8M\). Markov's inequality and \(b\le M\) imply
\[
\mathbb P(\|S\|>9M+v)\le\tfrac12\mathbb P(\|S\|>v).
\]
Thus \(\mathbb P(\|S\|>9jM)\le2^{-j}\) for \(j\ge0\), and integration
over consecutive intervals of length \(9M\) gives
\[
\E\|S\|^2
\le81M^2\sum_{j\ge0}(2j+1)2^{-j}
=486M^2<24^2M^2.
\]
\end{proof}

We will repeatedly use a simple consequence of quotient duality.
If a bounded surjection \(U:A\to B\) satisfies
\(\|U^*b^*\|\ge D^{-1}\|b^*\|\), then
\begin{equation}\label{eq:quotient-lift}
\inf_{Ua=b}\|a\|\le D\|b\|.
\end{equation}
Indeed, \(A/\ker U\) is isomorphic to \(B\), and its quotient norm at
\(b\) is
\[
\sup_{\|U^*b^*\|\le1}|\langle b,b^*\rangle|.
\]
In each application below surjectivity has an explicit bounded right
inverse, so the quotient identification is immediate.

\begin{lemma}[Bounded moment lifts]\label{lem:lifts}
Suppose
\[
c_m(G^*)\le D_g,\qquad c_m(F)\le D_f,
\qquad D_g,D_f\ge1.
\]
Set
\begin{equation}\label{eq:lift-parameters}
a=\frac1{2LD_g\sqrt m},\qquad
b=\frac1{2LD_f\sqrt m}.
\end{equation}
Then:
\begin{enumerate}[label=\textup{(\roman*)}]
\item For \(x_1,\ldots,x_m\in B_G\), there are \(P_\eps\in B_G\) such that
\[
\E_\eps\eps_iP_\eps=ax_i\quad(1\le i\le m).
\]
\item For \(y_1,\ldots,y_m\in B_{F^*}\), there are \(Z_\delta\in B_{F^*}\)
such that
\[
\E_\delta\delta_jZ_\delta=by_j\quad(1\le j\le m).
\]
\end{enumerate}
The same conclusions, with the same constants, hold after replacing
\(G,F\) by finite \(L_2\) sums of their copies.
\end{lemma}
\begin{proof}
The moment map
\[
U:L_\infty(\{-1,1\}^m;G)\longrightarrow\ell_2^m(G),
\qquad
UP=(\E_\eps\eps_iP_\eps)_{i=1}^m,
\]
is onto: a right inverse is \((x_i)\mapsto\sum_i\eps_i x_i\).
Its adjoint sends \((z_i^*)\) to \(\sum_i\eps_i z_i^*\).
By Lemma~\ref{lem:kahane} and \eqref{eq:fixed-cotype},
\[
\|U^*(z_i^*)\|_{L_1(G^*)}
\ge(LD_g)^{-1}\left(\sum_i\|z_i^*\|^2\right)^{1/2}.
\]
Equation \eqref{eq:quotient-lift} bounds the infimal lift norm of
\((ax_i)\) by \(LD_ga\sqrt m=1/2\). A lift of norm at most \(1\)
therefore exists, proving (i).

For (ii), the map
\[
V:\ell_2^m(F)\longrightarrow L_1(\{-1,1\}^m;F),
\qquad
V(w_j)=\sum_j\delta_jw_j,
\]
satisfies
\(\|Vw\|_{L_1}\ge(LD_f)^{-1}\|w\|_{\ell_2^m(F)}\).
The functional
\[
Vw\longmapsto b\sum_j y_j(w_j)
\]
has norm at most \(b\sqrt m\,LD_f=1/2\) on its range.
Hahn--Banach extends it to \(L_1(F)\) with the same norm.
Representing the extension in \(L_\infty(F^*)\) gives \(Z\) and the
stated moments. The amplification assertion follows from
\eqref{eq:cotype-amplification} and finite-sum duality.
\end{proof}
